% ============================================================
% main.tex — Submission 281
% "One Flood Substrate, Six Decision Geometries:
%  A Diagnostic Benchmark for Spatial AI [Experiment]"
% SIGSPATIAL 2026, Experiment/Benchmark & Experience track (<=10pp)
%
% OVERLEAF: New Project -> Upload Project -> drag the zip.
%   Compiler: pdfLaTeX. Main document: main.tex.
%   acmart class is built into Overleaf; not vendored here.
%
% BLINDING: default is single-blind (author block visible), matching the
%   submission record. If the track is double-blind, add `anonymous` to the
%   \documentclass options below and the author block auto-hides.
%
% FORMAT: sigconf is the SIGSPATIAL default. Confirm against the 2026
%   Experiment-track CFP; if different, change the class option on the
%   next line only.
% ============================================================
\documentclass[sigconf,nonacm]{acmart}
% ^ `nonacm` removes the ACM copyright/DOI block for clean DRAFT compilation.
%   DO NOT submit with nonacm without checking the SIGSPATIAL 2026 CFP: ACM
%   venues typically want plain [sigconf] (or [sigconf,review]) for the
%   submitted PDF. Confirm the required options before any reviewer-facing
%   build, and do not alter the template to gain space.
%   For double-blind: add `anonymous`.

\usepackage{booktabs}
\usepackage{pifont}      % \ding{51} \ding{55} in the availability matrix
\usepackage{graphicx}
\usepackage{xcolor}

% --- Placeholder + symbol macros (single source of truth) -----------------
% Every UNMEASURED number is wrapped in \PH{} and rendered red so it cannot
% be mistaken for a real result. Delete \PH once all results are locked.
\providecommand{\PH}[1]{\textcolor{red}{#1}}
\providecommand{\Rsup}{S_{\text{sup}}}
\providecommand{\TRF}{\mathrm{TRF}}  % Task Residual Floor (ADR-042, formerly N_ceiling)

% --- Front matter ----------------------------------------------------------
\begin{document}

\title{One Flood Substrate, Six Decision Geometries: A Diagnostic Benchmark for Spatial AI [Experiment]}
% NOTE (camera-ready): the [Experiment] suffix is required by the CFP at
% submission and removed at camera-ready.

\author{Rudolph A. Martin}
\affiliation{%
  \institution{Next Shift Consulting LLC}
  \city{Madison}
  \state{Wisconsin}
  \country{USA}
}
\email{rudy@nextshiftconsulting.com}

% NAMING NOTE: the abstract introduces the benchmark as "GeoRSCT," which is
% also the name of the NeurIPS D&B submission (martin2026georsct). Before
% camera-ready, disclose the relationship explicitly (this paper = GeoRSCT
% instantiated on a flood substrate; cite martin2026georsct as the
% machinery/validation) so the two are not read as duplicate submissions.

\begin{abstract}
Spatial AI benchmarks often reduce flood modeling to one held-out score.
That is insufficient for crisis decision support. A flood dataset may
support prediction, ranking, clustering, transfer, relational propagation,
or allocation, and each task can fail for a different spatial reason.

We introduce GeoRSCT, a diagnostic benchmark for evaluating which
flood-decision tasks a spatial dataset can support. GeoRSCT uses a
multi-event flood substrate spanning multiple metropolitan regions, with
radar and satellite precipitation, terrain, land cover, drainage networks,
levee and urban infrastructure, insurance claims, citizen-reported impacts,
and per-feature provenance and field-status annotations. The benchmark
organizes evaluation around six spatial decision geometries: prediction,
ranking, clustering, transfer, relational propagation, and allocation.

GeoRSCT does not ask only which model scores best. It asks whether the
evidence is sufficient for the decision being made: whether rankings change
under infrastructure context, whether model errors cluster spatially,
whether models transfer across storms and cities, whether drainage or
adjacency assumptions hold across built infrastructure, and whether rainfall
timing gives enough lead time for resource allocation. The result is a
diagnostic account of what a flood dataset supports, what it does not, and
where the evidence is missing or provisional --- surfacing, before a model
is trusted, the points at which a crisis manager would be acting without
sufficient spatial support.
\end{abstract}

\begin{CCSXML}
<ccs2012>
<concept>
<concept_id>10002951.10003227.10003351</concept_id>
<concept_desc>Information systems~Data mining</concept_desc>
<concept_significance>500</concept_significance>
</concept>
<concept>
<concept_id>10010147.10010178</concept_id>
<concept_desc>Computing methodologies~Machine learning</concept_desc>
<concept_significance>300</concept_significance>
</concept>
</ccs2012>
\end{CCSXML}
\ccsdesc[500]{Information systems~Data mining}
\ccsdesc[300]{Computing methodologies~Machine learning}
% TODO: replace with the CCS concepts you select from the ACM CCS tool.

\keywords{Spatial AI, GeoAI, spatial machine learning, spatial data
  infrastructure, geospatial benchmarks, data contracts, reproducible
  benchmarks, flood decision support, evidence sufficiency, disaster response}

\maketitle

% --- Body ------------------------------------------------------------------
% ============================================================
% SECTION 1 — Introduction
% ============================================================
\section{Introduction}
\label{sec:intro}

Flood models produce per-location damage estimates. A crisis manager
uses those estimates for at least six distinct operations: predicting
damage, prioritizing response, identifying where errors concentrate,
transferring models across storms or cities, reasoning about drainage
and adjacency, and allocating resources against forecast lead time.
Each operation can fail for a different spatial reason: prediction can
be inflated by spatial leakage across validation folds; rankings can
shift when infrastructure context enters the evidence; errors can
cluster in neighborhoods the model has not learned; transfer can fail
across events even when within-event accuracy is adequate; adjacency
can be severed by levees or pumps; and temporal evidence can be absent
for events that predate a sensor network.

A single held-out accuracy score addresses prediction and is silent on
the other five. Existing geospatial benchmarks~\cite{yeh2021sustainbench,
klemmer2023satclip, agarwal2025pdfm, krechetova2025geobenchx} report one
metric per task. That metric answers which model ranks best on average;
it does not answer whether the evidence is sufficient for the specific
decision a manager needs to make.

This paper introduces GeoRSCT, a diagnostic benchmark organized around
six spatial \emph{decision geometries}---prediction, ranking, clustering,
transfer, relational propagation, and allocation---each paired with a
named spatial failure mode (\S\ref{sec:problem-geometries}). For each
geometry, the benchmark applies a pre-committed, leakage-controlled test
and reports a sufficiency verdict: \textsc{supported}, \textsc{provisional},
or \textsc{insufficient}. The full set of verdicts is the benchmark's
output: a diagnostic account of what a flood substrate supports and where
the evidence is missing or inadequate.

The benchmark runs on a multi-event flood substrate spanning five U.S.\
metropolitan regions, with radar and satellite precipitation, terrain,
drainage networks, levee and urban infrastructure, insurance claims,
citizen-reported impacts, and high-water marks. Every feature carries
provenance and a field-status annotation so that absent data is
distinguished from a measured zero. The evaluation uses a strict
representation ladder (R0$\to$R1$\to$R2) with fixed folds, solvers, and
seeds; only the feature set changes between levels, so changes in verdict
are attributable to the added evidence. At each level, an RSCT
certificate~\cite{martin2026rsct} decomposes model quality into signal,
suppression, and noise components, providing an auditable record of
\emph{why} a verdict changed, not only that it did. The certificate is
computed and recorded before the next level trains, so the quality
assessment of level $k$ is independent of level $k{+}1$.

GeoRSCT is a benchmark, not a deployed system. It does not propose a
model or optimize a score. Its contribution is a reusable substrate,
a six-geometry evaluation protocol, and a worked example---including
negative results and evidence gaps---of what evidence sufficiency means
in practice for spatial decision support.

% ============================================================
% SECTION 2 — Related Work
% ============================================================
\section{Related Work}
\label{sec:related}

\paragraph{Spatial validation and leakage.}
Standard cross-validation assumes independent observations. When data
exhibit spatial autocorrelation, nearby locations leak information
across folds and inflate reported skill. Roberts et
al.~\cite{roberts2017cv} formalized spatial blocking strategies;
Ploton et al.~\cite{ploton2020spatial} demonstrated that large-scale
ecological mapping models can lose most of their apparent skill under
spatial validation; and Mil\`{a} et al.~\cite{mila2022nndm} proposed
nearest-neighbour distance matching for map-aware fold construction.
These methods address the prediction geometry. Our benchmark adopts
spatial blocking as the headline evaluation protocol and reports the
gap between random and blocked skill as the leakage diagnostic, but
extends the evaluation to five additional decision geometries that
spatial blocking alone does not address.

\paragraph{Scale, aggregation, and administrative-unit pitfalls.}
The Modifiable Areal Unit Problem~\cite{openshaw1984maup} established
that statistical relationships change with the spatial unit of
aggregation. Resolution mismatch between labels and predictions
introduces systematic error in remote sensing~\cite{workman2022resolution}.
In health and policy research, ZIP codes and Census-defined ZCTAs are
routinely conflated despite material spatial
mismatches~\cite{krieger2002zip, chen2025crosswalk}. Our substrate
operates at the ZCTA grain and documents the ZIP-as-ZCTA approximation
as a known limitation (C-001, \S\ref{sec:substrate}). Geographically
weighted regression~\cite{brunsdon1996gwr} provides a lens for
diagnosing spatial non-stationarity in model coefficients, which we
apply in the clustering and prediction geometries
(\S\ref{sec:spatial-diagnostics}).

\paragraph{Geospatial benchmarks.}
SustainBench~\cite{yeh2021sustainbench} evaluates sustainability
indicators across multiple tasks but reports per-task accuracy, not
per-decision sufficiency. SatCLIP~\cite{klemmer2023satclip} benchmarks
location embeddings; PDFM~\cite{agarwal2025pdfm} evaluates a
geospatial foundation model on population dynamics tasks with
state-holdout extrapolation. GeoYCSB~\cite{kim2023geoycsb} targets
database performance, not model evaluation.
GeoBenchX~\cite{krechetova2025geobenchx} and
GeoAnalystBench~\cite{zhang2025geoanalyst} benchmark agent-based
geospatial workflows. Dynabench~\cite{kiela2021dynabench} and
Ott et al.~\cite{ott2022saturation} identify benchmark saturation
as a systemic problem in ML evaluation. Each of these contributes a
dataset, a task, or a platform; none organizes evaluation around the
decision the prediction is used for or reports per-geometry evidence
sufficiency. That is the gap this benchmark addresses.

\paragraph{RSCT and GeoRSCT.}
The Representation--Solver Compatibility Theory
(RSCT)~\cite{martin2026rsct} decomposes model quality into a simplex
of signal ($R$), suppression ($S_\text{sup}$), and noise ($N$), with
derived measures for signal purity ($\alpha$), geometric compatibility
($\kappa$), and cross-fold stability ($\tau$, $\sigma$). GeoRSCT
(NeurIPS D\&B)~\cite{martin2026georsct} validates this machinery on a
27-task CONUS-scale geospatial benchmark. The present paper instantiates
the same certificate framework on a flood substrate organized around
decision geometries rather than regression targets, using the
certificate as an audit trail that explains verdict changes across
representation levels.

\input{sections/03_problem_formulation}
% ============================================================
% SECTION 4 — The Flood Substrate
% ============================================================
\section{The Flood Substrate}
\label{sec:substrate}

\begin{figure*}[t]
  \centering
  \includegraphics[width=\textwidth]{figures/nfip_scenario_comparison.pdf}
  \caption{Cumulative NFIP claim footprint across the five CONUS benchmark
    scenarios, for scale and orientation. These are cumulative all-hazard
    claims (through 2024) within each scenario's geographic footprint
    ---\emph{not} the event-matched per-(\textit{zcta\_id},\textit{event})
    target the benchmark models (\S\ref{sec:problem-substrate}). Oahu is
    excluded here: it is an out-of-distribution probe
    (Appendix~\ref{app:oahu}), not a benchmark scenario.}
  \label{fig:nfip-scenarios}
\end{figure*}

The benchmark operates on a single multi-event flood substrate at the
(\textit{zcta\_id}, \textit{event}) grain, spanning five U.S.\ metropolitan
regions selected for complementary hazard profiles and data availability
(Table~\ref{tab:scenarios}). Houston provides the densest claim record
(396 ZCTA--event pairs) with both tropical and inland flooding; Southwest
Florida (606 pairs) covers hurricane surge and coastal inundation; New York
City (422 pairs) adds nor'easter and tidal flood events; Riverside--Coachella
Valley (172 pairs) contributes arid flash-flood events with thin insurance
penetration; and New Orleans (264 pairs) spans the pre- and post-levee
era, offering a natural experiment in infrastructure-mediated flood risk.

\begin{table}[t]
\centering
\caption{Benchmark scenarios. $n$ is ZCTA--event rows. Targets:
  \texttt{nfip} = NFIP event claims (regression); \texttt{311} = citizen
  reports (classification); \texttt{hwm} = USGS high-water marks
  (classification). A cell is modelable when the target has enough
  non-missing rows for spatial-blocked CV.}
\label{tab:scenarios}
\small
\begin{tabular}{@{}llcl@{}}
\toprule
\textbf{Scenario} & \textbf{$n$} & \textbf{Targets} & \textbf{Hazard profile} \\
\midrule
Houston         & 396 & nfip, 311, hwm & Tropical + inland \\
SW Florida      & 606 & nfip           & Hurricane surge \\
NYC             & 422 & nfip, 311, hwm & Nor'easter + tidal \\
Riverside       & 172 & nfip           & Flash flood (arid) \\
New Orleans     & 264 & nfip, hwm      & Pre/post-levee \\
\bottomrule
\end{tabular}
\end{table}

\subsection{Data layers}

Features are organized by representation level
(\S\ref{sec:method-ladder}), where each level is a strict superset of
the previous:

\paragraph{R0 (static tabular, 33--36 features).}
American Community Survey demographics~\cite{kwon2023acs}, CDC Social
Vulnerability Index, FEMA flood zone classifications, terrain (TWI,
slope, elevation from USGS 3DEP), HIFLD critical infrastructure counts,
and temporally-gated NFIP claim history (computed strictly before each
event's onset to prevent temporal leakage).

\paragraph{R1 (+hydrology and spatial structure, 58--63 features).}
R0 plus NHDPlus drainage network metrics, USACE levee proximity, estimated
sewershed capacity, and eight spatial-lag features derived from a
queen-contiguity adjacency matrix $\mathbf{W}$. The target spatial lag
(\textit{wlag\_nfip\_claims}) is recomputed per fold from training units
only to prevent test-fold leakage (\S\ref{app:relational-ablation}).

\paragraph{R2 (+temporal dynamics, 67--72 features).}
R1 plus NOAA MRMS Stage~IV quantitative precipitation estimates, HRRR
forecast fields, tide-gauge readings, and HURDAT2 storm-track proximity.
These sources have operational start dates (MRMS: October 2012; HRRR:
September 2014), so pre-dating events carry a
\texttt{not\_applicable\_temporal} field status rather than an imputed
value.

\subsection{Provenance and field status}

Each feature carries a provenance record (source agency, vintage,
spatial resolution) and one of four field-status values:
\texttt{available} (value present), \texttt{not\_applicable\_temporal}
(event predates the source's operational start),
\texttt{not\_applicable\_type} (source does not apply to this hazard
type), or \texttt{source\_empty} (source covers the event but returned
zero records). This vocabulary is the mechanism by which the benchmark
surfaces evidence gaps rather than silently imputing through them---the
allocation verdict, for instance, is necessarily partial because
temporal features carry \texttt{not\_applicable\_temporal} for pre-2012
events (Appendix~\ref{app:protocol-substrate},
Table~\ref{app:tab:availability}).

\subsection{Known limitations}

\paragraph{ZIP-as-ZCTA (C-001).} NFIP claims are reported by ZIP code;
ZCTAs are Census-defined approximations. Within the five metropolitan
footprints, fewer than 3\% of ZIPs map to a different ZCTA. We treat
ZIP as ZCTA and document the approximation, following standard
practice~\cite{krieger2002zip, chen2025crosswalk}. The mismatch rate is
higher in rural areas outside our study regions.

\paragraph{Riverside sample size.} At $n{=}172$ ZCTA--event pairs,
Riverside--Coachella Valley is the thinnest scenario and produces the
widest confidence intervals. Results from this scenario should be
interpreted with that constraint.

% ============================================================
% SECTION 5 — Benchmark Construction (Method)
% ============================================================
\section{Benchmark Construction}
\label{sec:method}

The benchmark is a controlled representation intervention on a fixed
flood substrate. The same folds, same solvers, same targets, and same
evaluation protocol are held constant across three representation
levels; only the feature set changes. This section describes what is
built (the ladder), how inference is disaggregated (encode then decode),
and how pre-training quality signals drive operational adaptation (the
gearbox).

% ------------------------------------------------------------------
\subsection{Encode then decode}
\label{sec:encode-decode}

Geospatial prediction is encode-then-decode. The \emph{encoder}
identifies and joins the geography: which polygon is which ZCTA, which
addresses fall in which boundary, which units are neighbors, which
crosswalk version is current. This is classification work---discrete,
combinatorial, and Fano-native---and it produces the feature substrate
that everything downstream consumes. The \emph{decoder} consumes that
substrate: regression to predict continuous targets (claims, presence,
high-water marks), then a decision layer that allocates resources under
operational constraints.

The RSCT certificate primitives split cleanly across the two halves:
\begin{itemize}
\item \textbf{Encoder side.} Signal purity $\alpha = R/(R+N)$ is
  grounded where Fano's inequality applies natively: boundary
  classification, crosswalk identification, and spatial join correctness
  are discrete channel problems. The geometric compatibility
  $\kappa_\text{geom}$ (Phase~0.5) is also encoder-side: it measures
  spatial connectivity, feature coverage, and scale stability from
  problem structure alone, before any model is trained.
\item \textbf{Decoder side.} The task residual floor ($\TRF$) estimator
  is the decoder-side analog of Fano's noise-floor role: a
  rate-distortion lower bound on achievable prediction error,
  regime-appropriate for continuous targets. Cross-validation fold
  metrics (R$^2$, RMSE, ROC-AUC) are decoder-side measurements.
\item \textbf{Decision layer.} The gate scaffolding---sequential
  gates producing \textsc{execute}, \textsc{caution}, or
  \textsc{refuse}---operates at the boundary between prediction and
  operational use, where the question shifts from ``how accurate?'' to
  ``accurate enough for this decision?''
\end{itemize}

This separation matters because the benchmark's central claim---that
audit flags predict which representation fixes help---requires
encoder-side signals ($\kappa_\text{geom}$, spatial structure) to
predict decoder-side outcomes (fold-level uplift). If both sides used
the same measurements, the prediction would be circular.

% ------------------------------------------------------------------
\subsection{The representation ladder}
\label{sec:ladder}

Three representation levels form a strict superset chain. No features
are removed or re-tuned between levels, so any performance change is
attributable to the added features alone.

\begin{description}
\item[R0 (static tabular).] ACS demographics, SVI vulnerability indices,
  FEMA flood zone percentages, terrain derivatives (TWI, slope), HIFLD
  infrastructure proximity, and temporally-gated NFIP claim history.
  Approximately 33--36 features depending on scenario.
\item[R1 (hydrology + spatial lag).] R0 plus NHDPlus catchment
  attributes, levee proximity, sewershed type (NYC), and eight W-matrix
  spatial lag features computed per fold from training data only to
  prevent target leakage (Appendix~\ref{app:relational-ablation}).
  Approximately 58--63 features.
\item[R2 (temporal event dynamics).] R1 plus MRMS precipitation timing,
  tide gauge anomaly, and HURDAT2 storm-track features. Approximately
  67--72 features. Temporal sources have vintage boundaries
  (Table~\ref{app:tab:availability}); events predating the source are
  marked \texttt{not\_applicable\_temporal}, not imputed.
\end{description}

Solvers are held constant across levels: HistGBDT
(\texttt{max\_depth=6}, \texttt{learning\_rate=0.1}, \texttt{max\_iter=200})
and Ridge ($\alpha{=}1.0$, median imputation, standard scaling). No
hyperparameter tuning is performed. Random seed is 42. Fold assignments
are generated once (spatial blocking by county or ZIP3 fallback) and
reused at all levels.

% ------------------------------------------------------------------
\subsection{The gearbox: pre-training quality signals}
\label{sec:gearbox}

The encode/decode separation creates a problem for operational routing:
$\kappa_\text{geom}$ is encoder-side and level-invariant---it measures
problem geometry, not model quality---so it cannot distinguish between
R0, R1, and R2 for a given cell. A routing system that uses
$\kappa_\text{geom}$ alone will always recommend the same arm regardless
of representation level.

The \emph{gearbox} (Phase~0.75) sits at the boundary between encode and
decode. It runs a cheap pre-training probe---Ridge cross-validation on
the R0 feature set---and derives five quality signals per
(scenario, target) cell:

\begin{description}
\item[$\tau$ (stability).] $\tau = 1/(1 + \text{CV})$ where
  $\text{CV} = \text{std}(\text{folds}) / |\text{mean}(\text{folds})|$.
  High $\tau$ means stable fold performance; low $\tau$ means the cell
  is sensitive to fold composition.
\item[$\sigma$ (turbulence).] $\sigma = \text{std}(\text{fold metrics},
  \text{ddof}{=}1)$. Raw fold variance, independent of mean level.
\item[$\alpha_\text{warmup}$ (signal strength).] For regression,
  $\alpha = \max(0, R^2)$; for classification,
  $\alpha = 2 \cdot (\text{AUC} - 0.5)$. Rescaled to $[0,1]$ so both
  task types share a common quality axis.
\item[Collapse risk.] Fraction of folds where the Ridge solver fails to
  beat the mean predictor ($R^2 \leq 0$ or AUC $\leq 0.5$).
\item[Coherence.] $1 - \text{IQR} / |\text{median}|$, measuring
  whether folds agree on the direction and scale of signal.
\end{description}

These five signals are then calibrated into a discrete \emph{gear
assignment} per cell, following the GearboxCalibrator methodology: tau
values are collected across all cells, winsorized, and partitioned into
percentile-based boundaries with monotonicity guards. Each cell receives
a gear from a five-level system (Precision, Balanced, Explore, Discovery,
Reverse) that determines how aggressively the decoder should operate.

The gearbox output feeds into Phase~6 (DGM routing), which evaluates
two strategies side by side:
\begin{enumerate}
\item \textbf{Kappa-first (legacy).} Uses $\kappa_\text{geom}$ as the
  primary discriminant. Known limitation: level-invariant, so it always
  exits on the first branch.
\item \textbf{Gear-informed (new).} Uses gearbox signals (gear,
  $\alpha_\text{warmup}$, $\sigma$) as the primary discriminant, falling
  back to certificate comparison only when signals are ambiguous. These
  signals vary per cell, enabling actual discrimination across arms.
\end{enumerate}

The gear-informed strategy produces internal decisions
(\textsc{execute}, \textsc{re\_encode}, \textsc{repair},
\textsc{fallback}) that are projected to public decisions
(\textsc{execute}, \textsc{caution}, \textsc{refuse}) at the output
boundary---two vocabularies deliberately separated so that system-internal
routing logic does not leak into user-facing operational guidance.

% ------------------------------------------------------------------
\subsection{Pre-registration and tamper evidence}
\label{sec:preregistration}

Diagnostics and $\kappa_\text{geom}$ are computed and uploaded to S3
with timestamps \emph{before} the next training level executes.
Phase~0.5 (geometry kappa) precedes Phase~1 (R0 training); Phase~0.75
(gearbox warmup) generates its own folds from spatial structure rather
than consuming folds produced by Phase~1, maintaining independence. S3
object timestamps provide tamper-evident ordering: the predictor
($\kappa_\text{geom}$, warmup signals) is frozen before the outcome
(R1/R2 uplift) exists.

The full pre-registered protocol---hypotheses, design matrix, controlled
variables, statistical tests, kill rules, and change control---is in
Appendix~\ref{app:doe}.

\input{sections/06_evaluation}   % <-- the one fully-drafted section
\input{sections/07_results}
% ============================================================
% SECTION 8 — Conclusion
% ============================================================
\section{Conclusion}
\label{sec:conclusion}

This paper introduced GeoRSCT, a diagnostic benchmark that evaluates
flood substrates against the decisions they are used for, not against a
single held-out accuracy score. By organizing evaluation around six
spatial decision geometries---prediction, ranking, clustering, transfer,
relational propagation, and allocation---the benchmark surfaces where a
substrate's evidence is sufficient, where it is provisional, and where
it is absent, before a model is trusted for operational use.

The findings are mixed by design. Prediction is the geometry with the
strongest evidential support: the representation ladder produces
measurable, leakage-controlled skill that improves with added structure.
Ranking and clustering are descriptive: the probes detect reordering
under infrastructure context and residual spatial autocorrelation, but
the verdicts rest on nine cells and carry the uncertainty appropriate to
that sample size. Transfer (leave-event-out retention) is reported with
its gaps. Relational propagation is specified but not yet executed---the
W-matrix ablation variants are designed, and the verdict cannot be issued
until they run. Allocation is necessarily partial: events predating the
temporal sources carry \texttt{not\_applicable\_temporal} annotations,
and the benchmark reports this as an evidence gap, not a defect.

Several limitations bound the findings. The benchmark operates on
$n{=}9$ modelable cells, enough for one pre-registered confirmatory
test but too few for robust cell-level associations (H2b, H4 are
explicitly exploratory). The Riverside scenario, at 172 ZCTA--event
pairs, is the thinnest cell and produces the widest confidence
intervals. The ZIP-as-ZCTA approximation (C-001) is standard practice
but introduces spatial mismatch, particularly in rural areas outside
our study footprints. And the substrate is U.S.-specific: NFIP claims,
FEMA flood zones, and ACS demographics do not transfer to non-U.S.\
contexts without equivalent data infrastructure.

The benchmark's contribution is not a model or a score but a protocol:
a reusable substrate with provenance and field-status annotations, a
six-geometry evaluation framework with pre-committed verdict criteria,
and a worked example---including negative results and evidence
gaps---of what it means to ask whether the data supports the decision,
not just whether the model fits the data.


% --- Bibliography ----------------------------------------------------------
% SIGSPATIAL: 10 pages excluding references; up to 2 appendix pages AFTER
% references. Bibliography therefore precedes the appendix.
\bibliographystyle{ACM-Reference-Format}
\bibliography{refs}

% --- Appendix (after references; <=2 pages) --------------------------------
\appendix
% ============================================================
% APPENDIX for Submission 281
% "One Flood Substrate, Six Decision Geometries:
%  A Diagnostic Benchmark for Spatial AI [Experiment]"
%
% Purpose: paper-supporting appendix. Substantiates the abstract's two
% load-bearing claims —
%   (1) "the benchmark organizes evaluation around six spatial decision
%        geometries"  -> the protocol (evidence mapping + verdict logic)
%   (2) "per-feature provenance and field-status annotations" + "where the
%        evidence is missing or provisional" -> the substrate (sources,
%        field-status semantics, availability matrix)
%
% Scope discipline: this appendix carries ONLY what the abstract names.
% No certificate gate, action vocabulary, gearbox, oobleck, MoE/forecast
% stack, or controlplane architecture appears here — none is in the
% abstract, so none belongs in an appendix to it.
%
% Integration:
%   - assumes \appendix has been declared in the body
%   - \PH{} is the red placeholder macro (defined defensively below if absent)
%   - requires \usepackage{booktabs} and \usepackage{pifont}
% ============================================================


\section{Decision-Geometry Evaluation Protocol and Substrate}
\label{app:protocol-substrate}

\paragraph{Two distinct ``sixes.''}
The benchmark is organized around six \emph{decision} geometries---the
operational task types a crisis manager performs: prediction, ranking,
clustering, transfer, relational propagation, and allocation. These are
\emph{not} the six measurement proxies of the underlying compatibility
machinery (smooth, spatial, composite, hierarchical, logical, periodic),
which are diagnostic instruments. The two are orthogonal: a kappa proxy
may inform several decision geometries, and no decision geometry maps
one-to-one onto a single proxy. Indeed, as
Table~\ref{app:tab:evidence-map} shows, five of the six decision
geometries are evaluated by \emph{direct} comparison of model outputs
(rank correlations, leave-event-out deltas, ablation uplift); only
\emph{clustering} is evaluated through a registered kappa proxy
(\textit{kappa\_spatial}, via Moran's $I$ on residuals). The evaluation
axis is the decision, not the proxy.

\subsection{Protocol: per-geometry evidence}
\label{app:protocol}
Table~\ref{app:tab:evidence-map} states, for each decision geometry, the
test that supplies its evidence, the diagnostic it relies on, and its
implementation status as of this writing. We report status honestly:
results are not locked, and one geometry's evidence path is specified but
not yet implemented in the training code.

\begin{table*}[t]
\centering
\caption{Decision-geometry evidence map. ``Test'' is the artifact that
  supplies evidence; ``Diagnostic'' is what the test relies on (only
  clustering uses a registered kappa proxy); ``Status'' is the
  implementation state, reported honestly pending data lock. Every
  reported result downstream of this table is a placeholder until runs
  complete.}
\label{app:tab:evidence-map}
\small
\begin{tabular}{@{}p{2.5cm}p{4.3cm}p{3.4cm}p{4.2cm}@{}}
\toprule
\textbf{Decision geometry} & \textbf{Test (evidence artifact)} & \textbf{Diagnostic} & \textbf{Status} \\
\midrule
Prediction &
  R0$\to$R1$\to$R2 representation ladder; fold-level Wilcoxon on paired deltas &
  ladder skill ($R^2$, RMSE) &
  Specified; R0 executing \\[3pt]
Ranking &
  Within-event rank shift R0 vs R1 (Kendall's $\tau$-b, top-$k$ churn, fidelity-to-truth $\Delta$) &
  none --- direct rank comparison &
  Specified (ranking-geometry design); not executed \\[3pt]
Clustering &
  Moran's $I$ on spatial-blocked residuals, per level &
  \textit{kappa\_spatial} (registered) &
  Wired (\texttt{compute\_diagnostics.py} imports \texttt{yrsn.core.kappa.spatial}); not executed \\[3pt]
Transfer &
  Leave-event-out vs.\ spatial-blocked metric ratio &
  none --- direct metric comparison &
  Specified in all training scripts; not executed \\[3pt]
Relational propagation &
  W-matrix ablation (R1-full vs.\ R1-no-wlag vs.\ no-target-lag vs.\ wlag-only) &
  none --- direct uplift comparison (\textit{kappa\_logical} is the conceptual instrument, not computed) &
  \textbf{Specified in DOE; not implemented.} Training script produces R1-full only; ablation-variant predictions pending \\[3pt]
Allocation &
  R2 temporal lead-time features (MRMS/HRRR rainfall timing) &
  none --- direct (\textit{kappa\_periodic} is the conceptual instrument, not computed) &
  \textbf{Partial.} Coverage gap by data vintage (\S\ref{app:substrate}, Table~\ref{app:tab:availability}) \\
\bottomrule
\end{tabular}
\end{table*}

\subsection{Per-geometry sufficiency verdicts}
\label{app:verdicts}
The benchmark's output is not a leaderboard but a \emph{sufficiency
verdict} per geometry: \textsc{supported}, \textsc{provisional}, or
\textsc{insufficient}. The verdict criteria are fixed before results.

\begin{description}
\item[Prediction.] \textsc{Supported} if the ladder beats the mean
  predictor and a random-feature null under spatial-blocked CV
  ($R^2>0$ on $\geq 2$ targets) and R1 improves on R0 at the fold level
  (Wilcoxon $p<0.05$ \emph{and} Cohen's $d>0.2$); else \textsc{provisional}.
\item[Ranking.] With reordering threshold pre-committed at Kendall's
  $\tau\text{-b}<0.90$: \textsc{supported} if context meaningfully reorders
  \emph{and} moves the order toward observed impact
  ($\Delta\text{-fidelity}>0$); \textsc{provisional} if it reorders without
  improving fidelity; context not load-bearing if $\tau\text{-b}\geq 0.90$.
  Because NFIP claims are zero-inflated, the headline is the non-zero
  (claimant) subset when tied pairs exceed half.
\item[Clustering.] \textsc{Supported} if residual Moran's $I$ is
  indistinguishable from zero after the spatial ladder level;
  \textsc{insufficient} if errors remain spatially clustered (a standing
  evidence gap for any spatially-resolved decision).
\item[Transfer.] \textsc{Supported} if leave-event-out skill retains a
  pre-committed fraction of spatial-blocked skill; \textsc{provisional}
  otherwise. Cross-event is reported; cross-\emph{city} is weaker evidence
  and not pooled.
\item[Relational propagation.] \textsc{Supported} if R0$\to$R1 uplift
  survives the no-target-lag ablation; if uplift is entirely the target
  spatial lag, the honest verdict is \textsc{provisional}---``neighbor
  activity predicts,'' not ``relational structure is supported.''
  \emph{This verdict cannot be issued until the ablation variants are
  implemented} (Table~\ref{app:tab:evidence-map}).
\item[Allocation.] \textsc{Supported} for events with temporal coverage;
  \textsc{insufficient (evidence absent by vintage)} for events predating
  the rainfall-timing sources (\S\ref{app:substrate}). The verdict is
  therefore necessarily \emph{partial} across the event panel---a
  demonstration of the benchmark surfacing where a decision lacks support,
  not a defect.
\end{description}

\subsection{Substrate, provenance, and field status}
\label{app:substrate}
The benchmark runs on one multi-event flood substrate at the
(\textit{zcta\_id}, \textit{event}) grain, spanning five metropolitan
regions. The nine modelable cells---the basis of the $n{=}9$ figure used
throughout---are the (scenario, target) pairs with sufficient coverage to
train, enumerated in Table~\ref{app:tab:cells}.

\begin{table}[t]
\centering
\caption{The nine modelable (scenario, target) cells. A cell is modelable
  when the target is populated at the (\textit{zcta\_id}, \textit{event})
  grain with enough non-missing rows to fit under spatial-blocked CV;
  $n$ is rows (ZCTAs $\times$ events). This table is what grounds the
  $n{=}9$ claim.}
\label{app:tab:cells}
\small
\begin{tabular}{@{}llll@{}}
\toprule
\textbf{Scenario} & \textbf{Target} & \textbf{Task} & \textbf{$n$ rows} \\
\midrule
Houston      & \texttt{obs\_nfip\_event\_claims} & regression     & 396 \\
Houston      & \texttt{obs\_has\_311}            & classification & 396 \\
Houston      & \texttt{obs\_has\_hwm}            & classification & 396 \\
SW Florida   & \texttt{obs\_nfip\_event\_claims} & regression     & 606 \\
NYC          & \texttt{obs\_nfip\_event\_claims} & regression     & 422 \\
NYC          & \texttt{obs\_has\_311}            & classification & 422 \\
Riverside    & \texttt{obs\_nfip\_event\_claims} & regression     & 172 \\
New Orleans  & \texttt{obs\_nfip\_event\_claims} & regression     & 264 \\
New Orleans  & \texttt{obs\_has\_hwm}            & classification & 264 \\
\bottomrule
\end{tabular}
\end{table}

Each feature carries provenance and a \emph{field-status}
annotation drawn from a fixed vocabulary, so that an absent value is never
silently treated as a measured zero:

\begin{description}
\item[\texttt{available}.] The source covers this event; a value is present.
\item[\texttt{not\_applicable\_temporal}.] The event predates the source's
  operational start. This is a vintage boundary, \emph{not} a missing
  measurement.
\item[\texttt{not\_applicable\_type}.] The source does not apply to this
  event type (e.g.\ Atlantic best-track for an atmospheric-river event).
\item[\texttt{source\_empty}.] The source covers the event but returned
  zero records.
\end{description}

This annotation is what makes the allocation verdict legible. Allocation
depends on rainfall-timing features (NOAA MRMS Stage~IV; HRRR) and soil
moisture (SMAP). MRMS~v12 is operational from October~2012, HRRR from
September~2014, and SMAP from March~2015. Every event in the panel falls
within coverage \emph{except} the two earliest New Orleans events, which
are marked \texttt{not\_applicable\_temporal} rather than missing
(Table~\ref{app:tab:availability}). For those events, temporal lead-time
sufficiency cannot be assessed, and allocation is reported as
\textsc{insufficient (evidence absent by vintage)}.

\begin{table}[t]
\centering
\caption{Temporal-source coverage for the allocation geometry, shown on
  the New Orleans event series (the pre/post-levee substrate, where the
  vintage boundary bites). \ding{51}\,\texttt{available};
  \ding{55}\,\texttt{not\_applicable\_temporal} (predates source). All
  non-NOLA panel events postdate October~2012 and carry MRMS.}
\label{app:tab:availability}
\small
\begin{tabular}{@{}llccc@{}}
\toprule
\textbf{Event} & \textbf{Date} & \textbf{MRMS} & \textbf{HRRR} & \textbf{SMAP} \\
\midrule
Katrina & 2005-08-29 & \ding{55} & \ding{55} & \ding{55} \\
Isaac   & 2012-08-28 & \ding{55} & \ding{55} & \ding{55} \\
Barry   & 2019-07-11 & \ding{51} & \ding{51} & \ding{51} \\
Ida     & 2021-08-26 & \ding{51} & \ding{51} & \ding{51} \\
\midrule
\multicolumn{2}{@{}l}{Source operational from} & 2012-10 & 2014-09 & 2015-03 \\
\bottomrule
\end{tabular}
\end{table}

\paragraph{Reproducibility note.}
The field-status vocabulary above is the single source of truth consumed
by the data-readiness validator; an event/source pair marked
\texttt{not\_applicable\_temporal} is excluded from the affected geometry's
evidence rather than imputed, so the ``missing or provisional'' verdicts
are determined by the substrate annotation, not by post-hoc judgment.

\subsection{Relational ablation and leakage controls}
\label{app:relational-ablation}
The relational-propagation geometry (\S\ref{sec:geom-relational}) is
evaluated by ablating the R1 W-matrix block. The W-matrix adds eight
spatial-structure features to the R0$\to$R1 step: node degree, mean
neighbor distance, four static neighbor lags, a rainfall lag
(\textit{wlag\_rainfall\_mm}), and the lag of the target itself
(\textit{wlag\_nfip\_claims}). Four feature sets isolate what the spatial
structure contributes:

\begin{table}[h]
\centering
\caption{W-matrix ablation battery. Feature counts are for the Houston
  cell; the design is identical across scenarios.}
\label{app:tab:ablation}
\small
\begin{tabular}{@{}llc@{}}
\toprule
\textbf{Variant} & \textbf{Features} & \textbf{$n_\text{feat}$} \\
\midrule
full          & R0 + hydrology + W-matrix              & 61 \\
no-wlag       & R0 + hydrology only                    & 53 \\
no-target-lag & full minus \textit{wlag\_nfip\_claims} & 60 \\
wlag-only     & R0 + W-matrix only                     & 41 \\
\bottomrule
\end{tabular}
\end{table}

\paragraph{Verdict logic.} The relational verdict rests on
\textbf{no-target-lag}: if the R0$\to$R1 uplift survives removing
\textit{wlag\_nfip\_claims}, spatial structure is supported as evidence;
if the uplift is carried entirely by the target lag, the honest finding is
``neighbor activity predicts,'' reported as provisional. no-wlag and
wlag-only bound the contributions of point features and hydrology
respectively.

\paragraph{Target-lag leakage and its control.}
\textit{wlag\_nfip\_claims} is the spatial lag of the prediction target.
Computed once over the full dataset, it carries test-fold target values
into training: under spatially-blocked cross-validation a unit's neighbors
may sit in the held-out fold, so a training unit's lag is built partly from
held-out labels. This inflates precisely the relational uplift the ablation
is meant to measure---the test of whether spatial structure is real would
itself be contaminated by target-lag leakage. We therefore recompute
\textit{wlag\_nfip\_claims} \emph{per fold} from training units only: for
each fold, a unit's lag is the row-standardized mean of its neighbors'
target values restricted to the training partition, left undefined when all
of a unit's neighbors fall in the test fold. The variants that contain the
column (full, wlag-only) require an adjacency graph and are refused at
run time if it is absent, rather than silently falling back to the leaked
pre-computed values; the variants that do not (no-wlag, and the
headline no-target-lag) are unaffected and run without adjacency. The
headline relational verdict is thus clean of target-lag leakage by
construction.

\paragraph{Residual caveats.} The non-target spatial lags
(\textit{wlag\_rainfall\_mm} and the static neighbor lags) remain
full-dataset quantities; as lags of \emph{inputs} rather than the target,
they carry the ordinary spatial-smoothing dependence common to all spatial
features, not answer leakage, and are not recomputed per fold. A
residual-based diagnostic (\textit{spatial\_lag\_residual\_R0}) requiring
out-of-fold R0 predictions is specified but deferred to the R0 results
lock.

% ============================================================
% APPENDIX B — Out-of-distribution clustering probe (Oahu, s036)
% Real measured numbers (NOT \PH). Framed as instrument-transfer, not a
% sixth benchmark scenario; honest inconclusive-negative, no "unbiased" claim.
% ============================================================
\section{Out-of-Distribution Probe: Clustering on a Non-CONUS Substrate (Oahu)}
\label{app:oahu}

To test whether the clustering instrument (\textit{kappa\_spatial},
\S\ref{app:protocol-substrate}) transfers off the CONUS substrate, we apply
it to a structurally different setting: Oahu, Hawaii. This is a separate
experiment from the main benchmark---a different model (the Floodcaster
HAZUS damage engine, not the s035 representation ladder), a different region
(island topology, no continental adjacency), and a single decision geometry
(clustering only). It is an instrument-transfer probe, \emph{not} a sixth
benchmark scenario, and it carries no claim about the other five geometries.

\begin{figure*}[t]
  \centering
  \includegraphics[width=\textwidth]{figures/oahu_h4_results.pdf}
  \caption{Out-of-distribution clustering probe on Oahu (experiment s036,
    Floodcaster/HAZUS engine). The H4 clustering hypothesis \emph{fails}:
    residual autocorrelation is negative (Moran's $I=-0.358$, dispersion
    not clustering; \textit{kappa\_spatial}$=0.352$). At $n{=}17$ ZCTAs,
    3.5\% inundation, and a cumulative (non-event-matched) NFIP reference,
    this is an \emph{inconclusive negative}---not evidence that the damage
    model is geographically unbiased. The finding is a scope condition on
    the instrument (see text).}
  \label{fig:oahu}
\end{figure*}

% Companion figure (Floodcaster HAZUS depth--damage engine).
% Generated by: render_hazus_depth_damage.py --upload
\begin{figure*}[t]
  \centering
  \includegraphics[width=\textwidth]{figures/hazus_depth_damage_curves.pdf}
  \caption{FEMA HAZUS depth--damage functions used by the Floodcaster
    engine for the Oahu probe (\emph{not} the benchmark's HistGBDT/Ridge
    model). Panel~C shows Oahu observed damage by depth band.}
  \label{fig:hazus}
\end{figure*}

\paragraph{Setup.} A 10-year coastal-surge inundation raster with reef
attenuation (\texttt{Oahu\_10\_withReef}) drives HAZUS damage estimates for
47{,}983 buildings; 1{,}676 (3.5\%) are inundated (median depth 1.74\,ft,
94\% single-family slab-on-grade). Predicted losses are aggregated to the 17
of 20 Hawaii ZCTAs that contain buildings, compared against \emph{cumulative}
NFIP historical claims, normalized to $[0,1]$, and the absolute residual is
taken. Global Moran's $I$ is computed on the 17-ZCTA queen-contiguity graph.

\begin{table}[h]
\centering
\caption{Oahu clustering probe vs.\ the CONUS substrate. The probe returns
  negative (dispersed), not positive (clustered), residual autocorrelation.}
\label{app:tab:oahu}
\small
\begin{tabular}{@{}ll@{}}
\toprule
\textbf{Quantity} & \textbf{Value} \\
\midrule
Buildings / inundated        & 47{,}983 / 1{,}676 (3.5\%) \\
ZCTAs modeled                & 17 (of 20 HI) \\
Moran's $I$ (residuals)      & $-0.358$ \\
Null expectation $E[I]$      & $-0.0625$ \\
\textit{kappa\_spatial}      & $0.352$ (below the 0.5 neutral point) \\
Reference                    & cumulative NFIP (not event-matched) \\
\bottomrule
\end{tabular}
\end{table}

\paragraph{Result and interpretation.} Residual autocorrelation is
\emph{negative} ($I=-0.358$ vs.\ null $-0.0625$): adjacent ZCTAs tend to have
\emph{dissimilar} error magnitudes, the opposite of clustering. No positive
geographic clustering of residuals is detected, so the hypothesis that
\textit{kappa\_spatial} would flag reef/coastal discontinuities is not
supported on this substrate. We do \emph{not} read this as evidence that the
damage model is geographically unbiased: the same dispersed value is the
expected signature of an \emph{underpowered or confounded} test. Four
conditions make the negative result inconclusive rather than reassuring:
(i) at 3.5\% inundation, most ZCTAs carry near-zero predicted loss, so the
few wet ZCTAs alternate with dry neighbors and induce anti-correlation by
construction; (ii) ZCTA aggregation cross-cuts the reef line rather than
aligning with it, so a real coastal discontinuity is averaged away;
(iii) the reference is cumulative NFIP history, not the surge event, so the
residual measures model-vs-history mismatch, not model error; and
(iv) $n{=}17$ ZCTAs gives Moran's $I$ low power, and the sign could be noise.

\paragraph{What the probe contributes.} The value is a \emph{scope condition
on the instrument}, not a verdict on the model: \textit{kappa\_spatial} is
diagnostic for the clustering geometry only when wet-cell density and an
event-matched reference are adequate; under extreme inundation sparsity and a
historical reference it returns dispersion that must not be interpreted as
absence of geographic bias. For the Oahu substrate the clustering geometry is
therefore reported as \textsc{evidence-insufficient}---a clean demonstration
that the diagnostic, like any instrument, has an operating envelope the
benchmark should declare.

% ============================================================
% APPENDIX D — Design of Experiments
% Pre-registered experimental protocol for the s035 model-ladder.
% Covers hypotheses, design matrix, controlled variables, statistical
% tests, kill rules, and change control.  Written before results.
% ============================================================
\section{Design of Experiments}
\label{app:doe}

This appendix documents the pre-registered experimental protocol for the
representation ladder (s035). The protocol was locked before Phase~1
execution and amended twice under change control (Table~\ref{app:tab:doe-versions}).
All hypotheses, statistical tests, and kill rules were specified before
any model was trained. Post-hoc adjustments to pass verification are
forbidden by the verification gates (Appendix~\ref{app:agent-lifecycle},
Table~\ref{app:tab:verify-lifecycle}).

\subsection{Hypotheses}
\label{app:doe-hypotheses}

Four hypotheses are tested across the representation ladder. H2a is the
single pre-registered confirmatory test; the others are descriptive or
exploratory.

\begin{table*}[t]
\centering
\caption{Pre-registered hypotheses. H2a is confirmatory (Wilcoxon + effect
  size); all others are descriptive or exploratory. Kill rules halt the
  experiment early if the baseline fails.}
\label{app:tab:hypotheses}
\small
\begin{tabular}{@{}p{0.8cm}p{4.5cm}p{3.5cm}p{2.5cm}p{3.0cm}@{}}
\toprule
\textbf{ID} & \textbf{Statement} & \textbf{Test} & \textbf{Pass criterion} & \textbf{Type} \\
\midrule
H1 &
  R0 features achieve skill above the mean predictor on $\geq 2$ of 3 targets &
  $R^2 > 0$ under spatial-blocked CV &
  $\geq 2$ targets with $R^2 > 0$ &
  Descriptive baseline gate \\[4pt]
H2a &
  R1 improves over R0 at the fold level &
  Wilcoxon signed-rank on paired (R0, R1) fold metrics &
  $p < 0.05$ AND Cohen's $d > 0.2$ &
  \textbf{Confirmatory} \\[4pt]
H2b &
  $\kappa_\text{geom}$ predicts which cells benefit more from R1 &
  Spearman $\rho$($\kappa_\text{geom}$, uplift) with bootstrap CI &
  $\rho > 0$, 95\% CI excludes zero &
  Exploratory ($n{=}9$) \\[4pt]
H3 &
  R2 temporal features improve over R1 when temporal mismatch is active &
  Paired delta: metric$_\text{R2}$ $-$ metric$_\text{R1}$ per fold &
  Uplift $> 3\%$ on primary target &
  Descriptive \\[4pt]
H4 &
  Audit flags predict representation uplift &
  Spearman $\rho$(flag count, uplift) &
  $\rho > 0.3$, bootstrap CI excludes zero &
  Exploratory \\
\bottomrule
\end{tabular}
\end{table*}

\subsection{Design matrix}
\label{app:doe-design}

The experiment is a controlled representation intervention: same folds,
same solver, same target---only the feature set changes. This is causal
at the pipeline level (representation is the sole intervention) but not
at the hydrology level (we do not randomize watersheds).

\begin{table}[t]
\centering
\caption{Design matrix. The only factor that changes across arms is the
  representation level (R0/R1/R2). All other factors are held constant
  or crossed.}
\label{app:tab:design-matrix}
\small
\begin{tabular}{@{}llc@{}}
\toprule
\textbf{Factor} & \textbf{Levels} & \textbf{Type} \\
\midrule
Scenario        & 5 metros (Table~\ref{app:tab:cells}) & Observational \\
Target          & 3 (claims, 311, HWM)                  & Fixed \\
Representation  & R0, R1, R2                            & \textbf{Intervention} \\
Solver          & HistGBDT, Ridge                       & Fixed \\
Split protocol  & random, spatial-blocked, leave-event-out & Fixed \\
Fold            & 5-fold spatial-blocked CV             & Fixed \\
\bottomrule
\end{tabular}
\end{table}

\paragraph{Modelable cells.}
Not all (scenario, target) pairs are modelable: \texttt{obs\_has\_311}
requires a 311 system (Houston, NYC only); \texttt{obs\_has\_hwm}
requires high-water marks (Houston, New~Orleans only). The nine modelable
cells are enumerated in Table~\ref{app:tab:cells}.

\subsection{Representation levels}
\label{app:doe-repr}

Each level is a strict superset of the prior, ensuring that any
performance change is attributable to the added features, not to
dropped features or re-tuning.

\begin{table}[t]
\centering
\caption{Representation ladder. Each level adds features to the prior;
  no features are removed or re-tuned.}
\label{app:tab:repr-levels}
\small
\begin{tabular}{@{}lp{4.5cm}c@{}}
\toprule
\textbf{Level} & \textbf{Added features} & \textbf{$\sim n_\text{feat}$} \\
\midrule
R0 & Static tabular: ACS demographics, SVI, FEMA flood zones,
     terrain (TWI, slope), HIFLD infrastructure,
     temporally-gated NFIP history & 33--36 \\
R1 & R0 + hydrology (NHDPlus catchments, levee proximity,
     sewershed type) + W-matrix spatial lags (8 cols) & 58--63 \\
R2 & R1 + temporal event dynamics (MRMS precipitation timing,
     tide gauge anomaly, HURDAT2 track features) & 67--72 \\
\bottomrule
\end{tabular}
\end{table}

\paragraph{Feature count ranges.}
Approximate counts vary by 1--3 features across scenarios due to
scenario-specific columns (e.g., HCFCD drainage district for Houston,
sewershed type for NYC). Runtime validation is authoritative; documented
counts are approximate.

\subsection{Controlled variables}
\label{app:doe-controlled}

The following are held constant across all representation levels to
isolate the representation intervention:

\begin{table}[t]
\centering
\caption{Controlled variables. No hyperparameter tuning is performed;
  defaults are fixed before execution.}
\label{app:tab:controlled}
\small
\begin{tabular}{@{}lll@{}}
\toprule
\textbf{Variable} & \textbf{Value} & \textbf{Rationale} \\
\midrule
HistGBDT max\_iter      & 200     & No tuning \\
HistGBDT max\_depth     & 6       & No tuning \\
HistGBDT learning\_rate & 0.1     & No tuning \\
Ridge alpha             & 1.0     & No tuning \\
Random seed             & 42      & Reproducibility \\
Imputation (Ridge)      & median  & No tuning \\
Scaling (Ridge)         & StandardScaler & No tuning \\
Fold assignments        & Fixed, saved to S3 & Shared across R0/R1/R2 \\
\bottomrule
\end{tabular}
\end{table}

\subsection{Statistical tests}
\label{app:doe-tests}

\paragraph{H2a (confirmatory).}
Fold-level Wilcoxon signed-rank test on paired deltas
(metric$_\text{R1}$ $-$ metric$_\text{R0}$) across all folds and cells.
The test requires \emph{both} $p < 0.05$ and Cohen's $d > 0.2$; a small
$p$-value without meaningful effect size does not pass.

\paragraph{H2b (exploratory).}
Spearman rank correlation between $\kappa_\text{geom}$ (pre-training
geometric compatibility, computed from problem structure alone) and
per-cell uplift (R1 $-$ R0). Bootstrap 95\% CI ($B = 1000$) must exclude
zero. At $n{=}9$ cells this is explicitly exploratory.

\paragraph{All other tests.}
H1 is a descriptive gate (does skill exist?). H3 and H4 are descriptive
with bootstrap confidence intervals. No multiplicity correction is
applied because only H2a is confirmatory; the others are reported as
exploratory effect sizes with uncertainty.

\subsection{Kill rules}
\label{app:doe-kill}

Kill rules halt the experiment early and prevent wasted compute:

\begin{enumerate}
\item \textbf{H1 FAIL on all 3 targets.} Data quality is too low for any
  modeling. Stop. Do not proceed to R1.
\item \textbf{H2a: Wilcoxon $p \geq 0.05$ or Cohen's $d < 0.2$.}
  Report as null result (R1 does not improve on R0). Proceed to R2 but
  flag as negative evidence.
\item \textbf{All uplift $< 1\%$.} Representation differences are noise;
  R0 is sufficient. Report and stop.
\end{enumerate}

\subsection{DO NOT constraints}
\label{app:doe-donot}

These constraints are binding for the duration of the experiment:
\begin{itemize}
\item Do NOT change fold assignments between representation levels.
\item Do NOT tune hyperparameters per representation level.
\item Do NOT run on scenarios without completed data assembly.
\item Do NOT use GPU instances (tabular models only).
\item Do NOT add features to R0 after the baseline is locked.
\item Do NOT impute targets (drop rows with NaN targets).
\end{itemize}

\subsection{Four-document experiment governance}
\label{app:doe-governance}

Each experiment is governed by four complementary documents with
non-overlapping scope:

\begin{table}[t]
\centering
\caption{Four-document governance stack. Each document answers a different
  question; together they form the complete experiment record.}
\label{app:tab:four-docs}
\small
\begin{tabular}{@{}lp{5.0cm}@{}}
\toprule
\textbf{Document} & \textbf{Question answered} \\
\midrule
DOE (this appendix) &
  \emph{Why} are we running this and \emph{what} do we expect to find?
  Hypotheses, design matrix, statistical tests, kill rules. \\[4pt]
Experiment Contract &
  \emph{What} inputs and outputs does each phase need?
  Packages, S3 artifacts, feature schemas, verification gate definitions. \\[4pt]
Checklist &
  \emph{Where} are we in execution?
  Checkboxes, dates, job names, commit hashes. \\[4pt]
Verification Gates &
  \emph{How} do we know results are correct?
  V0--V5 gates with commands, pass criteria, failure protocol. \\
\bottomrule
\end{tabular}
\end{table}

The DOE is pre-registration: it prevents post-hoc hypothesis invention.
The contract ensures data dependencies are met before launch. The
checklist tracks progress. The verification gates (Appendix~\ref{app:agent-lifecycle})
ensure correctness. Amendments to the DOE are permitted under change
control (Table~\ref{app:tab:doe-versions}) but must be recorded before
the affected phase executes.

\subsection{Change control}
\label{app:doe-changes}

\begin{table}[t]
\centering
\caption{DOE version history. Amendments are permitted before the affected
  phase executes; post-hoc amendments are forbidden.}
\label{app:tab:doe-versions}
\small
\begin{tabular}{@{}lll@{}}
\toprule
\textbf{Version} & \textbf{Date} & \textbf{Change} \\
\midrule
v1.0 & 2026-05-29 & Initial DOE (draft) \\
v1.1 & 2026-05-29 & Locked for execution \\
v1.2 & 2026-06-01 & R1/R2 feature bundles replaced with verified S3 sources;
                     kappa proxy diagnostics added \\
v1.7 & 2026-06-03 & H2 redesigned: fold-level Wilcoxon + Cohen's $d$ \\
v1.8 & 2026-06-03 & H2b added: $\kappa_\text{geom}$ independence test \\
\bottomrule
\end{tabular}
\end{table}

\paragraph{Paper-safe framing.}
This experiment is causal at the pipeline level: same folds, same solver,
same target---only the representation changes. It is \emph{not} causal at
the hydrology level (we do not randomize watersheds). The correct framing
is: ``Under controlled representation intervention, scenarios flagged by
audit B.1 showed $X$\% larger improvement when hydrologic features were
added.'' The incorrect framing---``adding hydrologic features causes better
flood prediction''---is confounded by feature quality, spatial coverage,
and data vintage.

% ============================================================
% APPENDIX E — Notation
% Summary of mathematical notation used throughout the paper.
% ============================================================
\section{Notation}
\label{app:notation}

Table~\ref{app:tab:notation} collects the notation used throughout the
paper. Symbols are grouped by category: substrate and observation,
decision geometries, representation ladder, certificate primitives,
statistical tests, and spatial diagnostics.

\begin{table}[t]
\centering
\caption{Notation summary.}
\label{app:tab:notation}
\small
\begin{tabular}{@{}lp{5.8cm}@{}}
\toprule
\textbf{Symbol} & \textbf{Definition} \\
\midrule
\multicolumn{2}{@{}l}{\textit{Substrate and observation}} \\[2pt]
$\mathcal{S}$ & Flood substrate tuple $(\mathcal{L}, \mathcal{E}, \mathcal{X}, \mathcal{Y}, \mathcal{P})$ \\
$\mathcal{L}$ & Set of locations (ZCTAs) \\
$\mathcal{E}$ & Set of flood events \\
$\mathcal{X}$ & Feature map over $\mathcal{L} \times \mathcal{E}$ \\
$\mathcal{Y}$ & Target map over $\mathcal{L} \times \mathcal{E}$ \\
$\mathcal{P}$ & Per-feature provenance and field-status annotation \\
$(\ell, e)$ & Observation unit: one ZCTA $\times$ one event \\
$\hat{y}_{\ell,e}$ & Model prediction for location $\ell$, event $e$ \\
$n$ & Number of ZCTA--event rows in a scenario \\
\midrule
\multicolumn{2}{@{}l}{\textit{Decision geometries and verdicts}} \\[2pt]
$g$ & A decision geometry (one of six) \\
$v_g(\mathcal{S})$ & Sufficiency verdict for geometry $g$ on substrate $\mathcal{S}$: \textsc{supported}, \textsc{provisional}, or \textsc{insufficient} \\
\midrule
\multicolumn{2}{@{}l}{\textit{Representation ladder}} \\[2pt]
R0 & Static tabular features ($\sim$33--36 cols) \\
R1 & R0 + hydrology + spatial lag ($\sim$58--63 cols) \\
R2 & R1 + temporal event dynamics ($\sim$67--72 cols) \\
$\mathbf{W}$ & Queen-contiguity adjacency matrix \\
\midrule
\multicolumn{2}{@{}l}{\textit{Certificate primitives}} \\[2pt]
$R$ & Signal component (simplex; $R + S_\text{sup} + N = 1$) \\
$S_\text{sup}$ & Suppression component \\
$N$ & Noise component \\
$\alpha$ & Signal purity: $R / (R + N)$ \\
$\kappa_\text{geom}$ & Geometric compatibility (encoder-side, pre-training) \\
$\kappa_\text{compat}$ & Compatibility proxy: $R \cdot (1 - N)$ \\
$\tau$ & Fold stability: $1 / (1 + \text{CV})$ \\
$\sigma$ & Cross-fold turbulence: $\text{std}(\text{fold metrics})$ \\
$\TRF$ & Task Residual Floor (rate-distortion lower bound) \\
\midrule
\multicolumn{2}{@{}l}{\textit{Statistical tests}} \\[2pt]
$d$ & Cohen's $d$ (standardized effect size) \\
$\tau_b$ & Kendall's $\tau$-b (rank correlation) \\
$\rho$ & Spearman rank correlation \\
$\Delta$ & Fidelity change: $\tau_b(\text{R1}, \text{truth}) - \tau_b(\text{R0}, \text{truth})$ \\
\midrule
\multicolumn{2}{@{}l}{\textit{Spatial diagnostics}} \\[2pt]
$I$ & Moran's $I$ (global spatial autocorrelation) \\
$E[I]$ & Expected value of $I$ under spatial randomness \\
\bottomrule
\end{tabular}
\end{table}


\end{document}
